\documentclass[12pt,a4paper]{article}

% ==================== PACKAGES ====================
\usepackage[utf8]{inputenc}
\usepackage[vietnamese]{babel}
\usepackage{amsmath,amssymb,amsthm}
\usepackage{geometry}
\usepackage{enumitem}
\usepackage[most]{tcolorbox}
\usepackage{hyperref}
\usepackage{fancyhdr}
\usepackage{graphicx}
\usepackage{tikz}
\usepackage{eso-pic}
\usepackage{xcolor}

\geometry{a4paper, margin=2cm, bottom=2.5cm}
\hypersetup{colorlinks=true, linkcolor=blue!70!black, urlcolor=blue}
\setlist[itemize]{topsep=4pt, itemsep=2pt}
\setlist[enumerate]{topsep=4pt, itemsep=2pt}

% ==================== WATERMARK ====================
\AddToShipoutPictureFG{%
  \AtPageCenter{%
    \begin{tikzpicture}[remember picture, overlay]
      \node[rotate=45, scale=1, opacity=0.06]{%
        \fontsize{60}{72}\selectfont\bfseries\sffamily Đặng Tấn Quí%
      };
    \end{tikzpicture}%
  }%
}

% ==================== HEADER / FOOTER ====================
\pagestyle{fancy}
\fancyhf{}
\fancyhead[L]{\small\textit{Tích phân và Sự thay đổi tích lũy}}
\fancyhead[R]{\small\textit{Đặng Tấn Quí}}
\fancyfoot[C]{\thepage}
\fancyfoot[L]{\footnotesize\textcopyright\ Đặng Tấn Quí}
\fancyfoot[R]{\footnotesize SĐT: 0377\,122\,210}
\renewcommand{\headrulewidth}{0.4pt}
\renewcommand{\footrulewidth}{0.4pt}

\fancypagestyle{plain}{%
  \fancyhf{}
  \fancyfoot[C]{\thepage}
  \fancyfoot[L]{\footnotesize\textcopyright\ Đặng Tấn Quí}
  \fancyfoot[R]{\footnotesize SĐT: 0377\,122\,210}
  \renewcommand{\headrulewidth}{0pt}
  \renewcommand{\footrulewidth}{0.4pt}
}

% ==================== CUSTOM ENVIRONMENTS ====================
\newtcolorbox{dinhnghia}[1][]{enhanced, breakable, colback=blue!3!white, colframe=blue!60!black, fonttitle=\bfseries, title={Định nghĩa}, #1}
\newtcolorbox{dinhly}[1][]{enhanced, breakable, colback=green!3!white, colframe=green!50!black, fonttitle=\bfseries, title={Định lý}, #1}
\newtcolorbox{tinhchat}[1][]{enhanced, breakable, colback=yellow!5!white, colframe=orange!50!black, fonttitle=\bfseries, title={Tính chất}, #1}
\newtcolorbox{phuongphap}[1][]{enhanced, breakable, colback=gray!5!white, colframe=gray!50!black, fonttitle=\bfseries, title={Phương pháp}, #1}
\newtcolorbox{luuy}[1][]{enhanced, breakable, colback=red!3!white, colframe=red!50!black, fonttitle=\bfseries, title={Lưu ý}, #1}
\newtcolorbox{congthuc}[1][]{enhanced, breakable, colback=cyan!3!white, colframe=cyan!50!black, fonttitle=\bfseries, title={Công thức}, #1}
\newtcolorbox{vidu}[1][]{enhanced, breakable, colback=violet!3!white, colframe=violet!50!black, fonttitle=\bfseries, title={Ví dụ}, #1}

% ==================== DOCUMENT ====================
\begin{document}

% ==================== TRANG BÌA ====================
\thispagestyle{empty}
\begin{tikzpicture}[remember picture, overlay]
  \fill[blue!80!black] (current page.south west) rectangle (current page.north east);
  \fill[blue!60!cyan, opacity=0.8]
    (current page.south west) -- (current page.north west) --
    ([xshift=-3cm]current page.north east) -- (current page.south east) -- cycle;
  \foreach \r/\o in {6/0.05, 4.5/0.08, 3/0.12} {
    \fill[white, opacity=\o] ([xshift=2cm, yshift=-2cm]current page.north east) circle (\r cm);
  }
  \draw[white, line width=2pt] ([yshift=-5cm]current page.north west) ++(2cm,0) -- ++(17cm,0);
  \draw[white, line width=2pt] ([yshift=-16cm]current page.north west) ++(2cm,0) -- ++(17cm,0);
  \node[anchor=west, white, font=\fontsize{28}{34}\bfseries\selectfont]
    at ([xshift=3cm, yshift=-7.5cm]current page.north west)
    {PHẦN 4: TÍCH PHÂN};
  \node[anchor=west, white, font=\fontsize{20}{26}\bfseries\selectfont]
    at ([xshift=3cm, yshift=-9.2cm]current page.north west)
    {VÀ SỰ THAY ĐỔI TÍCH LŨY};
  \node[anchor=west, white, font=\fontsize{16}{20}\selectfont]
    at ([xshift=3cm, yshift=-11cm]current page.north west)
    {Tài liệu luyện thi du học};
  \node[white, opacity=0.12, font=\fontsize{120}{120}\selectfont, rotate=-15]
    at ([xshift=-3cm, yshift=4cm]current page.center)
    {$\displaystyle\int$};
  \node[anchor=west, white, font=\Large\bfseries]
    at ([xshift=3cm, yshift=-13cm]current page.north west)
    {Biên soạn:};
  \node[anchor=west, yellow!80!white, font=\fontsize{28}{34}\bfseries\selectfont]
    at ([xshift=3cm, yshift=-14.5cm]current page.north west)
    {ĐẶNG TẤN QUÍ};
  \node[anchor=west, white!90, font=\large]
    at ([xshift=3cm, yshift=-18cm]current page.north west)
    {SĐT: 0377\,122\,210};
  \node[anchor=south, white!80, font=\small]
    at ([yshift=1.5cm]current page.south)
    {Tài liệu có bản quyền --- Cấm sao chép, phân phối khi chưa được sự đồng ý của tác giả};
  \node[anchor=south east, white, font=\Large\bfseries]
    at ([xshift=-2cm, yshift=3cm]current page.south east) {\the\year};
\end{tikzpicture}
\newpage

\thispagestyle{empty}
\vspace*{5cm}
\begin{center}
  {\Large\bfseries PHẦN 4: TÍCH PHÂN VÀ SỰ THAY ĐỔI TÍCH LŨY}\\[8pt]
  {\large Tài liệu luyện thi du học}
\end{center}
\vspace{2cm}
\begin{tcolorbox}[enhanced, colback=red!3!white, colframe=red!60!black, fonttitle=\bfseries, title={Thông tin bản quyền}, boxrule=1.5pt]
  \textbf{Tác giả:} Đặng Tấn Quí \quad \textbf{Liên hệ:} 0377\,122\,210 \quad \textbf{Năm:} \the\year \\[6pt]
  \textbf{Bản quyền \textcopyright\ \the\year\ Đặng Tấn Quí. Bảo lưu mọi quyền.}
\end{tcolorbox}
\vfill\begin{center}\small\textit{Tài liệu lưu hành nội bộ}\end{center}
\newpage
\tableofcontents
\newpage

%% ================================================================
%%  PHẦN 4: TÍCH PHÂN VÀ SỰ THAY ĐỔI TÍCH LŨY
%% ================================================================
\section{Tích phân và Sự thay đổi tích lũy}

%% ----------------------------------------------------------------
%%  4.1 Tổng Riemann và Quy tắc hình thang
%% ----------------------------------------------------------------
\subsection{Tổng Riemann (Trái, Phải, Điểm giữa) và Quy tắc hình thang}

\begin{dinhnghia}[title={Tổng Riemann}]
  Để xấp xỉ diện tích dưới đường cong $y = f(x)$ trên $[a, b]$, chia đoạn thành $n$ khoảng bằng nhau, mỗi khoảng rộng $\Delta x = \dfrac{b-a}{n}$. Gọi $x_i = a + i\Delta x$ ($i = 0, 1, \ldots, n$). Khi đó:
  \begin{itemize}
    \item \textbf{Tổng Riemann trái:} $L_n = \Delta x \displaystyle\sum_{i=0}^{n-1} f(x_i) = \Delta x[f(x_0) + f(x_1) + \cdots + f(x_{n-1})]$.
    \item \textbf{Tổng Riemann phải:} $R_n = \Delta x \displaystyle\sum_{i=1}^{n} f(x_i) = \Delta x[f(x_1) + f(x_2) + \cdots + f(x_n)]$.
    \item \textbf{Tổng Riemann điểm giữa:} $M_n = \Delta x \displaystyle\sum_{i=1}^{n} f\!\left(\frac{x_{i-1}+x_i}{2}\right)$.
  \end{itemize}
\end{dinhnghia}

\begin{congthuc}[title={Quy tắc hình thang (Trapezoid Rule)}]
  \[
    T_n = \frac{\Delta x}{2}\bigl[f(x_0) + 2f(x_1) + 2f(x_2) + \cdots + 2f(x_{n-1}) + f(x_n)\bigr].
  \]
  Quy tắc hình thang chính xác hơn tổng Riemann trái/phải vì nó lấy trung bình: $T_n = \dfrac{L_n + R_n}{2}$.
\end{congthuc}

\begin{luuy}
  \begin{itemize}
    \item Nếu $f$ tăng trên $[a,b]$: $L_n \le \text{giá trị thực} \le R_n$ (đánh giá dưới và trên).
    \item Nếu $f$ giảm trên $[a,b]$: $R_n \le \text{giá trị thực} \le L_n$.
    \item Tổng Riemann là ý tưởng nền tảng cho định nghĩa tích phân xác định.
  \end{itemize}
\end{luuy}

\begin{vidu}
  Xấp xỉ $\displaystyle\int_0^2 x^2\,dx$ bằng $n = 4$ hình chữ nhật trái, phải, và quy tắc hình thang.

  \medskip\textbf{Giải:}
  $\Delta x = 0.5$. Các điểm: $x_0=0, x_1=0.5, x_2=1, x_3=1.5, x_4=2$.

  $f(x) = x^2$: $f(0)=0, f(0.5)=0.25, f(1)=1, f(1.5)=2.25, f(2)=4$.

  $L_4 = 0.5(0 + 0.25 + 1 + 2.25) = 0.5 \times 3.5 = 1.75$.

  $R_4 = 0.5(0.25 + 1 + 2.25 + 4) = 0.5 \times 7.5 = 3.75$.

  $T_4 = \dfrac{1.75 + 3.75}{2} = 2.75$.

  Giá trị thực: $\displaystyle\int_0^2 x^2\,dx = \dfrac{8}{3} \approx 2.667$.
\end{vidu}

%% ----------------------------------------------------------------
%%  4.2 Định lý cơ bản của Giải tích (FTC)
%% ----------------------------------------------------------------
\subsection{Định lý cơ bản của Giải tích (Fundamental Theorem of Calculus -- FTC)}

\begin{dinhly}[title={FTC -- Phần 1: Đạo hàm của hàm tích phân có cận biến}]
  Nếu $f$ liên tục trên $[a, b]$ và $g(x) = \displaystyle\int_a^x f(t)\,dt$, thì $g$ khả vi trên $(a, b)$ và:
  \[
    g'(x) = \frac{d}{dx}\int_a^x f(t)\,dt = f(x).
  \]
  \textbf{Dạng tổng quát (Chain Rule kết hợp):}
  \[
    \frac{d}{dx}\int_a^{u(x)} f(t)\,dt = f(u(x)) \cdot u'(x).
  \]
\end{dinhly}

\begin{dinhly}[title={FTC -- Phần 2: Công thức Newton--Leibniz}]
  Nếu $f$ liên tục trên $[a, b]$ và $F$ là nguyên hàm bất kỳ của $f$ ($F' = f$), thì:
  \[
    \int_a^b f(x)\,dx = F(b) - F(a) = \bigl[F(x)\bigr]_a^b.
  \]
\end{dinhly}

\begin{luuy}
  \begin{itemize}
    \item FTC nối liền \textbf{đạo hàm và tích phân}: hai phép toán ngược nhau.
    \item FTC--1 nói rằng tích phân là nguyên hàm (tích phân có cận biến là nguyên hàm của hàm dưới dấu tích phân).
    \item FTC--2 (Newton--Leibniz) cho phép tính tích phân xác định mà không cần tính tổng Riemann.
  \end{itemize}
\end{luuy}

\begin{vidu}
  Tính $\dfrac{d}{dx}\displaystyle\int_1^{x^3} \sin(t^2)\,dt$.

  \medskip\textbf{Giải:}
  Dùng FTC--1 kết hợp Chain Rule với $u = x^3$:
  \[
    \frac{d}{dx}\int_1^{x^3} \sin(t^2)\,dt = \sin\!\left[(x^3)^2\right] \cdot 3x^2 = 3x^2 \sin(x^6).
  \]
\end{vidu}

%% ----------------------------------------------------------------
%%  4.3 Tích phân bất định và nguyên hàm cơ bản
%% ----------------------------------------------------------------
\subsection{Tích phân bất định và các công thức nguyên hàm cơ bản}

\begin{dinhnghia}
  \[
    \int f(x)\,dx = F(x) + C \;\;\Leftrightarrow\;\; F'(x) = f(x),\quad C \in \mathbb{R}.
  \]
\end{dinhnghia}

\begin{congthuc}[title={Bảng nguyên hàm cơ bản}]
  \begin{align*}
    &\int x^n\,dx = \frac{x^{n+1}}{n+1} + C \quad (n \ne -1)
    &\quad&
    \int \frac{1}{x}\,dx = \ln|x| + C \\[5pt]
    &\int e^x\,dx = e^x + C
    &&
    \int a^x\,dx = \frac{a^x}{\ln a} + C \\[5pt]
    &\int \sin x\,dx = -\cos x + C
    &&
    \int \cos x\,dx = \sin x + C \\[5pt]
    &\int \tan x\,dx = \dfrac{1}{\cos^2 x} + C
    &&
    \int \cot x\,dx = \dfrac{-1}{\sin^2 x} + C \\[5pt]
    &\int \frac{1}{\sqrt{1-x^2}}\,dx = \arcsin x + C
    &&
    \int \frac{1}{1+x^2}\,dx = \arctan x + C \\[5pt]
    &\int \frac{dx}{x^2+a^2} = \frac{1}{a}\arctan\frac{x}{a} + C
    &&
    \int \frac{dx}{\sqrt{a^2-x^2}} = \arcsin\frac{x}{a} + C
  \end{align*}
\end{congthuc}

\begin{tinhchat}
  \begin{itemize}
    \item $\displaystyle\int [f(x) \pm g(x)]\,dx = \int f(x)\,dx \pm \int g(x)\,dx$.
    \item $\displaystyle\int c\cdot f(x)\,dx = c \int f(x)\,dx$ ($c$ là hằng số).
  \end{itemize}
\end{tinhchat}

%% ----------------------------------------------------------------
%%  4.4 U-Substitution
%% ----------------------------------------------------------------
\subsection{Phương pháp đổi biến (U-Substitution)}

\begin{phuongphap}[title={U-Substitution}]
  Dùng khi nhận ra biểu thức tích phân có dạng $\displaystyle\int f(g(x)) g'(x)\,dx$ (Chain Rule ngược):
  \begin{enumerate}
    \item Chọn $u = g(x)$ (thường là biểu thức ``bên trong'', phức tạp hơn).
    \item Tính $du = g'(x)\,dx$.
    \item Thế vào: $\displaystyle\int f(u)\,du = F(u) + C$.
    \item Thế $u$ trở lại theo $x$.
  \end{enumerate}
  \textbf{Với tích phân xác định:} đổi cận theo $u$ (không cần đổi $u$ về $x$ ở cuối).
\end{phuongphap}

\begin{vidu}
  Tính các tích phân sau:
  \begin{enumerate}[label=(\alph*)]
    \item $\displaystyle\int 2x\,e^{x^2}\,dx$
    \item $\displaystyle\int_0^1 x\sqrt{1+x^2}\,dx$
    \item $\displaystyle\int \frac{\cos(\ln x)}{x}\,dx$
  \end{enumerate}

  \medskip\textbf{Giải:}
  \begin{enumerate}[label=(\alph*)]
    \item Đặt $u = x^2$, $du = 2x\,dx$:
      $\displaystyle\int e^u\,du = e^u + C = e^{x^2} + C$.
    \item Đặt $u = 1+x^2$, $du = 2x\,dx$. Khi $x=0$: $u=1$; $x=1$: $u=2$.
      $\dfrac{1}{2}\displaystyle\int_1^2 \sqrt{u}\,du = \dfrac{1}{2}\cdot\dfrac{2}{3}u^{3/2}\Big|_1^2 = \dfrac{1}{3}(2^{3/2} - 1) = \dfrac{2\sqrt{2}-1}{3}$.
    \item Đặt $u = \ln x$, $du = \dfrac{1}{x}\,dx$:
      $\displaystyle\int \cos u\,du = \sin u + C = \sin(\ln x) + C$.
  \end{enumerate}
\end{vidu}

%% ----------------------------------------------------------------
%%  4.5 Tích phân từng phần
%% ----------------------------------------------------------------
\subsection{Tích phân từng phần (Integration by Parts)}

\begin{congthuc}[title={Công thức tích phân từng phần}]
  \[
    \int u\,dv = uv - \int v\,du.
  \]
  Dạng tích phân xác định:
  \[
    \int_a^b u\,dv = \bigl[uv\bigr]_a^b - \int_a^b v\,du.
  \]
\end{congthuc}

\begin{phuongphap}[title={Chọn $u$ và $dv$ -- Quy tắc LIATE}]
  Ưu tiên chọn $u$ theo thứ tự \textbf{LIATE}:
  \begin{enumerate}
    \item \textbf{L}ogarit: $\ln x$, $\log x$, \ldots
    \item \textbf{I}nverse trig: $\arcsin x$, $\arctan x$, \ldots
    \item \textbf{A}lgebraic: $x^n$, đa thức.
    \item \textbf{T}rig: $\sin x$, $\cos x$, \ldots
    \item \textbf{E}xponential: $e^x$, $a^x$, \ldots
  \end{enumerate}
  Phần còn lại (không chọn làm $u$) gán là $dv$.
\end{phuongphap}

\begin{vidu}
  Tính:
  \begin{enumerate}[label=(\alph*)]
    \item $\displaystyle\int x e^x\,dx$
    \item $\displaystyle\int x^2 \sin x\,dx$
    \item $\displaystyle\int \ln x\,dx$
  \end{enumerate}

  \medskip\textbf{Giải:}
  \begin{enumerate}[label=(\alph*)]
    \item Chọn $u = x$, $dv = e^x\,dx$ $\Rightarrow$ $du = dx$, $v = e^x$:
      $\displaystyle\int x e^x\,dx = xe^x - \int e^x\,dx = xe^x - e^x + C = (x-1)e^x + C$.

    \item Dùng từng phần hai lần. Chọn $u = x^2$, $dv = \sin x\,dx$:
      \begin{align*}
        \int x^2 \sin x\,dx &= -x^2 \cos x + 2\int x \cos x\,dx \\
        &= -x^2\cos x + 2\left(x\sin x - \int\sin x\,dx\right) \\
        &= -x^2\cos x + 2x\sin x + 2\cos x + C.
      \end{align*}

    \item Chọn $u = \ln x$, $dv = dx$ $\Rightarrow$ $du = \dfrac{1}{x}dx$, $v = x$:
      $\displaystyle\int \ln x\,dx = x\ln x - \int 1\,dx = x\ln x - x + C$.
  \end{enumerate}
\end{vidu}

%% ----------------------------------------------------------------
%%  4.6 Phân thức hữu tỉ
%% ----------------------------------------------------------------
\subsection{Tích phân bằng phân thức hữu tỉ (Partial Fractions)}

\begin{phuongphap}[title={Phân tích thành phân thức hữu tỉ (Partial Fraction Decomposition)}]
  Dùng để tính $\displaystyle\int \frac{p(x)}{q(x)}\,dx$ khi bậc $\deg(p) < \deg(q)$ và $q(x)$ phân tích được thành nhân tử.

  \textbf{Các dạng phân thức:}
  \begin{enumerate}
    \item \textbf{Mẫu có nhân tử bậc nhất phân biệt:} $\dfrac{A}{x-a} + \dfrac{B}{x-b} + \cdots$
    \item \textbf{Mẫu có nhân tử bậc nhất lặp lại:} $(x-a)^n \to \dfrac{A_1}{x-a} + \dfrac{A_2}{(x-a)^2} + \cdots + \dfrac{A_n}{(x-a)^n}$
    \item \textbf{Mẫu có nhân tử bậc hai bất khả quy:} $x^2 + bx + c \to \dfrac{Ax + B}{x^2 + bx + c}$
  \end{enumerate}
  \textbf{Nếu bậc tử $\ge$ bậc mẫu:} Chia đa thức trước, sau đó phân tích phần dư.
\end{phuongphap}

\begin{vidu}
  Tính $\displaystyle\int \frac{2x+1}{x^2 - x - 2}\,dx$.

  \medskip\textbf{Giải:}
  Phân tích: $x^2 - x - 2 = (x-2)(x+1)$.

  $\dfrac{2x+1}{(x-2)(x+1)} = \dfrac{A}{x-2} + \dfrac{B}{x+1}$.

  Nhân hai vế với $(x-2)(x+1)$: $2x+1 = A(x+1) + B(x-2)$.

  Cho $x = 2$: $5 = 3A \Rightarrow A = \dfrac{5}{3}$.
  Cho $x = -1$: $-1 = -3B \Rightarrow B = \dfrac{1}{3}$.

  \[
    \int \frac{2x+1}{x^2-x-2}\,dx = \frac{5}{3}\ln|x-2| + \frac{1}{3}\ln|x+1| + C.
  \]
\end{vidu}

\medskip\noindent\textbf{Các dạng bài tập:}
\begin{enumerate}[label=\textbf{Dạng \arabic*:}, leftmargin=*, align=left]
  \item Phân tích phân thức với mẫu có nhân tử thực phân biệt.
  \item Phân tích với mẫu có nhân tử lặp lại.
  \item Phân tích với mẫu có nhân tử bậc hai bất khả quy.
  \item Chia đa thức trước khi phân tích.
\end{enumerate}

%% ----------------------------------------------------------------
%%  4.7 Tích phân suy rộng
%% ----------------------------------------------------------------
\subsection{Tích phân suy rộng (Improper Integrals)}

\begin{dinhnghia}[title={Tích phân suy rộng loại 1 (Cận vô cực)}]
  \[
    \int_a^{+\infty} f(x)\,dx = \lim_{b \to +\infty} \int_a^b f(x)\,dx.
  \]
  Nếu giới hạn tồn tại (hữu hạn), tích phân \textbf{hội tụ}; ngược lại \textbf{phân kỳ}.

  Tương tự:
  \[
    \int_{-\infty}^b f(x)\,dx = \lim_{a \to -\infty} \int_a^b f(x)\,dx.
  \]
  \[
    \int_{-\infty}^{+\infty} f(x)\,dx = \int_{-\infty}^c f(x)\,dx + \int_c^{+\infty} f(x)\,dx \quad \text{(cả hai phải hội tụ)}.
  \]
\end{dinhnghia}

\begin{dinhnghia}[title={Tích phân suy rộng loại 2 (Hàm bị gián đoạn)}]
  Nếu $f$ không liên tục tại $b$:
  \[
    \int_a^b f(x)\,dx = \lim_{t \to b^-} \int_a^t f(x)\,dx.
  \]
  Nếu $f$ không liên tục tại $a$:
  \[
    \int_a^b f(x)\,dx = \lim_{t \to a^+} \int_t^b f(x)\,dx.
  \]
\end{dinhnghia}

\begin{dinhly}[title={Tích phân $p$}]
  \[
    \int_1^{+\infty} \frac{1}{x^p}\,dx \begin{cases} = \dfrac{1}{p-1} & \text{nếu } p > 1 \quad \text{(hội tụ)} \\ \text{phân kỳ} & \text{nếu } p \le 1 \end{cases}
  \]
\end{dinhly}

\begin{phuongphap}[title={Kiểm tra hội tụ bằng so sánh}]
  \textbf{Tiêu chuẩn so sánh:} Nếu $0 \le f(x) \le g(x)$ với $x \ge a$:
  \begin{itemize}
    \item $\displaystyle\int_a^{+\infty} g(x)\,dx$ hội tụ $\Rightarrow$ $\displaystyle\int_a^{+\infty} f(x)\,dx$ hội tụ.
    \item $\displaystyle\int_a^{+\infty} f(x)\,dx$ phân kỳ $\Rightarrow$ $\displaystyle\int_a^{+\infty} g(x)\,dx$ phân kỳ.
  \end{itemize}
  \textbf{Tiêu chuẩn so sánh giới hạn:} Nếu $\displaystyle\lim_{x\to\infty} \frac{f(x)}{g(x)} = L > 0$ (hữu hạn), thì $\int f$ và $\int g$ cùng hội tụ hoặc cùng phân kỳ.
\end{phuongphap}

\begin{vidu}
  Tính $\displaystyle\int_1^{+\infty} \frac{1}{x^2}\,dx$ và $\displaystyle\int_0^1 \frac{1}{\sqrt{x}}\,dx$.

  \medskip\textbf{Giải:}
  \begin{enumerate}[label=(\alph*)]
    \item $\displaystyle\int_1^{+\infty} \frac{1}{x^2}\,dx = \lim_{b\to\infty}\left[-\frac{1}{x}\right]_1^b = \lim_{b\to\infty}\left(-\frac{1}{b}+1\right) = 1$. Hội tụ.
    \item $\displaystyle\int_0^1 \frac{1}{\sqrt{x}}\,dx = \lim_{t\to 0^+}\left[2\sqrt{x}\right]_t^1 = \lim_{t\to 0^+}(2 - 2\sqrt{t}) = 2$. Hội tụ.
  \end{enumerate}
\end{vidu}

\medskip\noindent\textbf{Các dạng bài tập:}
\begin{enumerate}[label=\textbf{Dạng \arabic*:}, leftmargin=*, align=left]
  \item Tính tích phân suy rộng loại 1 (cận $+\infty$, $-\infty$).
  \item Tính tích phân suy rộng loại 2 (hàm gián đoạn trong đoạn).
  \item Kiểm tra hội tụ/phân kỳ bằng tiêu chuẩn so sánh.
  \item Kiểm tra hội tụ của tích phân $p$.
\end{enumerate}

\end{document}
